\chapter{Glosario ampliado}
\label{app:glosario}

\begin{longtable}{p{0.22\textwidth}p{0.66\textwidth}}
  \caption{Glosario de términos usados en la memoria.}
  \label{tab:glosario}\\
  \toprule
  Término & Definición \\
  \midrule
  \endfirsthead
  \toprule
  Término & Definición \\
  \midrule
  \endhead
  Alice & Nodo emisor en un protocolo QKD prepare-and-measure. \\
  Bob & Nodo receptor que mide los estados enviados por Alice. \\
  BB84 & Protocolo de distribución cuántica de claves propuesto por Bennett y Brassard en 1984. \\
  Base Z & Base computacional; en \textit{time-bin}, suele asociarse a llegada temprana o tardía. \\
  Base X & Base conjugada; en \textit{time-bin}, se mide mediante interferencia entre modos temporales. \\
  Canal clásico & Canal público autenticado usado para bases, tamizado, reconciliación y control. \\
  Canal cuántico & Medio óptico por el que viajan los estados cuánticos o pulsos débiles. \\
  Cascade & Protocolo interactivo de corrección de errores basado en comparación de paridades. \\
  Conteo oscuro & Click de detector no asociado a un fotón de señal. \\
  Decoy state & Pulso con intensidad distinta a la señal usado para estimar cotas de seguridad con fuentes coherentes débiles. \\
  Detector APD & Fotodiodo de avalancha usado para detección de fotones, común en implementaciones prácticas. \\
  Detector SNSPD & Detector superconductor de nanohilo, de alto rendimiento pero mayor complejidad operacional. \\
  Eve & Adversario o espía en la descripción criptográfica del protocolo. \\
  Fibra oscura & Fibra disponible sin tráfico clásico activo, útil para aislar el canal cuántico. \\
  Interferómetro Mach--Zehnder & Dispositivo óptico que separa y recombina caminos para medir superposiciones de fase o tiempo. \\
  Jitter & Incertidumbre temporal en emisión, propagación, detección o registro de eventos. \\
  KMS & Sistema de administración de claves que almacena, controla y entrega material de clave a aplicaciones. \\
  PNS & Photon Number Splitting; ataque contra pulsos multifotónicos de fuentes coherentes débiles. \\
  PQC & Criptografía post-cuántica basada en algoritmos clásicos resistentes a ataques cuánticos conocidos. \\
  QBER & Quantum Bit Error Rate; fracción de bits erróneos en la clave tamizada. \\
  QKD & Quantum Key Distribution; distribución cuántica de claves. \\
  SeQUeNCe & Simulador de eventos discretos para redes cuánticas usado en el proyecto. \\
  SKR & Secret Key Rate; tasa estimada de llave secreta luego de considerar postprocesamiento. \\
  Tamizado & Etapa en la que Alice y Bob conservan posiciones con bases compatibles. \\
  Time-bin & Codificación que representa información en modos temporales tempranos y tardíos. \\
  Trusted relay & Nodo intermedio confiable que participa en una red QKD de varios enlaces. \\
  Visibilidad & Contraste interferométrico; en el modelo se relaciona con error de fase. \\
  WCS & Weak Coherent Source; fuente coherente débil, típicamente un láser atenuado. \\
  \bottomrule
\end{longtable}

Este glosario está pensado como apoyo de lectura. En una versión final con defensa oral, puede convertirse también en una diapositiva de referencia para separar términos físicos, criptográficos y de red.
